\section{Conclusion}
\label{sec:conclusion}

This research shows how machine learning can be integrated into a unified, physics-based lunar-forecasting system while retaining the physical foundations of established environmental models. The development process answered the first research question by showing that a hazard-specific hybrid approach was appropriate for the current model, rather than applying machine learning uniformly across all hazards. The operational evaluation answered the second by showing that integrating hazards and data sources can improve how information is presented during mission planning, but it did not establish that LIPAS is safer than existing approaches. A unified system can also concentrate risk if its data or service fails. LIPAS is best considered alongside current NASA and NOAA systems as a consultative decision-support tool, subject to validation with real mission workflows and measurements.

The quantitative evaluation answered the third selectively. Chronological holdout RMSE was about 4.97 versus about 13.66 for persistence, and solar activity improved modestly in both windows. Radiation did not: hybrid RMSE exceeded the replica in night-dominant Window~A (0.119 versus 0.102) and in quiet dayside Window~B (0.0188 versus 0.00266). Temperature, moonquakes, and dust remained physics ties. Micrometeorite skill is circular under blend weight~0 and remains shower-term descriptive only. These studies therefore teach where the hybrid helps and where it still fails.

The central result is that a carefully constrained hybrid architecture can add predictive capability where machine learning is useful without replacing physics-based methods where they are more appropriate. LIPAS's principal contribution is therefore not the replacement of existing lunar models, but their integration into a single forecasting and decision-support framework.

The findings must be read with those limits in mind: proxy labels, quiet-regime radiation shortfalls, and circular micrometeorite scoring. LIPAS should function alongside NASA and NOAA products as consultative decision support, not as a replacement. The studies show how much remains to develop---radiation under quiet and storm regimes, stronger validation with surface measurements, and clearer operational workflows. This work does not end here; it is a first constrained step that the next iterations must extend.
